\documentclass[12pt, a4paper]{article}

% Packages pour le français et la mise en page
\usepackage[utf8]{inputenc}
\usepackage[T1]{fontenc}
\usepackage[french]{babel}
\usepackage[margin=2.5cm]{geometry}
\usepackage{parskip} % Pour aérer les paragraphes

\begin{document}

\begin{center}
    \LARGE \textbf{Notes d'oral : De la Renaissance au Siècle des Lumières}
\end{center}
\vspace{1cm}

\vspace{0.5cm}

\section*{PARTIE 1 : Comprendre l'Univers et la Nature}
Tout commence par une révolution scientifique majeure. Quelques génies décident que pour comprendre le monde, il ne faut plus lire les textes anciens, mais \textbf{observer} par soi-même.

\begin{itemize}
    \item Le pionnier absolu de cette démarche, c'est \textbf{Léonard de Vinci} (1452-1519). Symbole de la Renaissance, il refuse de croire aveuglément. Il observe la nature, étudie la botanique, et surtout, il pratique la dissection de cadavres pour comprendre la véritable anatomie humaine. Il incarne la curiosité pure.\\
    
    \item Le deuxième grand bouleversement vient de l'astronomie avec \textbf{Galilée} (1564-1642). Grâce à l'invention de la lunette astronomique, il observe le ciel et prouve l'\textbf{héliocentrisme} : c'est la Terre qui tourne autour du Soleil, et non l'inverse. C'est un choc monumental qui contredit la Bible. L'Homme n'est plus au centre de la création.\\
    
    \item Ensuite, \textbf{Johannes Kepler} (1571-1630) va plus loin. Grâce aux mathématiques, il prouve que les planètes ne tournent pas en cercles parfaits, mais forment des orbites en forme d'ellipses. L'univers n'est donc pas "magique", il obéit à la géométrie.\\
    
    \item Enfin, le sommet de cette révolution est atteint par l'Anglais \textbf{Isaac Newton} (1642-1727). Il découvre la \textbf{loi de la gravitation universelle}. Newton prouve qu'une seule et même force explique à la fois pourquoi une pomme tombe d'un arbre et pourquoi la Lune tourne autour de la Terre. \\
\end{itemize}

Pourquoi ces savants sont-ils si importants pour notre sujet ? Parce que de De Vinci à Newton, ils ont prouvé que la nature fonctionne comme une immense mécanique, avec des règles logiques que l'intelligence humaine peut comprendre. Cette idée va fasciner la génération suivante : les philosophes des Lumières.


\vspace{0.5cm}
\clearpage
\section*{PARTIE 2 : Éclairer la Société.}

Nous voici au 18e siècle, le fameux "Siècle des Lumières". Les intellectuels de cette époque se disent : \textit{"Si la Raison a permis de comprendre les lois de la nature, alors utilisons cette même Raison pour comprendre et améliorer les lois de la société !"} Leur but ? Chasser l'obscurantisme, c'est-à-dire l'ignorance et les préjugés. 

\begin{itemize}
    \item Pour cela, il faut éduquer. C'est le projet fou de \textbf{Denis Diderot} avec la création de l'\textbf{Encyclopédie}. Son but est de rassembler absolument toutes les connaissances (scientifiques, techniques, philosophiques) et de les diffuser au plus grand nombre.\\
    
    \item Ensuite, il y a \textbf{Voltaire}. Il est le grand défenseur de la liberté d'expression et de la tolérance. Inspiré par les sciences (il était un grand admirateur de Newton), il se bat avec sa plume contre le fanatisme religieux et les injustices politiques de son temps.\\
    
    \item D'autres vont s'attaquer à la politique pure. \textbf{Montesquieu} analyse les différents systèmes politiques de manière presque scientifique. Il comprend que le pouvoir absolu corrompt toujours. Il invente donc la \textbf{séparation des pouvoirs} (exécutif, législatif, judiciaire). Aujourd'hui, toutes nos démocraties modernes reposent sur cette idée.\\
    
    \item Enfin, \textbf{Jean-Jacques Rousseau} va s'attaquer aux inégalités. Avec son livre \textit{Le Contrat Social}, il détruit l'idée que le roi a un "droit divin" pour régner. Pour Rousseau, c'est le peuple qui est souverain. Les hommes naissent libres et égaux, et c'est la société mal construite qui les rend inégaux.\\
\end{itemize}

\section*{CONCLUSION.}

Enfin, comment résumer cette immense aventure intellectuelle qui a duré plus de deux siècles ? 

C'est le philosophe allemand \textbf{Emmanuel Kant} qui en donne la meilleure définition en 1784. Pour lui, les Lumières, c'est le moment où l'Humanité devient enfin adulte. Il résume cette pensée par une formule célèbre en latin : \textbf{"Sapere aude !"}, ce qui signifie \textit{"Aie le courage de te servir de ta propre intelligence"} (ou Ose penser par toi-même).

Les savants de la Renaissance ont osé observer le monde malgré les menaces de l'Église. Les philosophes des Lumières ont osé remettre en question le pouvoir absolu du roi. Ensemble, ils ont préparé le terrain à un événement qui allait faire exploser l'ancien monde de manière définitive, quelques années plus tard à peine : la \textbf{Révolution française de 1789}.


\end{document}